\section{Agregación en \texorpdfstring{\campo{resultados}}{resultados} y sensibilidad}
\label{sec:agregacion}

Esta sección cubre \campo{agregacion.py}, \campo{exportar.py} y
\campo{sensibilidad.py}, en \campo{src/analisis/}. Parten de los edificios
ya asignados (sección anterior) y calculan, para cada municipio, fuente y
umbral de confianza, las métricas de la ecuación~(1) de la Entrega~1.

\subsection{Colección \texorpdfstring{\campo{resultados}}{resultados}}

\campo{agregacion.py} crea \campo{resultados} con validación
\campo{\$jsonSchema} estricta, derivada del diccionario de la Entrega~1. Hay
un cambio respecto a ese diccionario: la llave pasa de
\campo{\{cod\_mpio, fuente\_id\}} a \campo{\{cod\_mpio, fuente\_id, umbral\}},
para guardar la sensibilidad en la misma colección sin crear otra. Microsoft
no tiene confianza y lleva \campo{umbral: null}. \campo{area\_media\_m2}
acepta \campo{null} porque un municipio sin edificios no tiene media.
\campo{parametros} guarda el umbral, la regla de asignación, el CRS de las
áreas y el commit del código que produjo el documento.

\subsection{Pipeline de agregación}

Cada combinación de fuente y umbral es una sola pipeline que corre dentro de
MongoDB:

\begin{lstlisting}[language=Python]
[{"$match": {"cod_mpio": {"$ne": None},
             "confianza": {"$gte": u}}},     # solo Google
 {"$group": {"_id": "$cod_mpio", "n": {"$sum": 1},
             "area": {"$sum": "$area_m2"},
             "media": {"$avg": "$area_m2"}}},
 {"$lookup": {"from": "municipios", "localField": "_id",
              "foreignField": "_id", "as": "m"}},
 {"$unwind": "$m"},
 {"$project": {"_id": {"cod_mpio": "$_id", "fuente_id": f, "umbral": u},
   "n_edificios": "$n", "area_total_m2": "$area", "area_media_m2": "$media",
   "densidad_km2": {"$divide": ["$n", "$m.area_km2"]},
   "pct_techos": {"$multiply": [100, {"$divide":
       ["$area", {"$multiply": ["$m.area_km2", 1e6]}]}]},
   "parametros": {...}, "fecha_calculo": "$$NOW"}},
 {"$merge": {"into": "resultados", "on": "_id",
             "whenMatched": "replace"}}]
\end{lstlisting}

Una segunda pipeline, sobre \campo{municipios}, completa con ceros los
municipios que no tienen edificios (\campo{\$merge} con
\campo{whenMatched: "keepExisting"}), así siempre hay 170 documentos por
fuente y umbral. Antes de escribir se borran los documentos de esa fuente y
umbral, de modo que correr el script dos veces deja exactamente lo mismo: en
la segunda corrida se reemplazaron los 170 documentos de cada combinación y
las cifras no cambiaron. Los umbrales se leen de \campo{config.yaml}
(\campo{analisis.umbrales\_google: [0.65, 0.70, 0.75]}), no están escritos en
el código. Cada pipeline tardó entre 4,2 y 4,5\,s en Google y 1,9\,s en
Microsoft, y cada corrida queda en \campo{bitacora\_cargas}
(\campo{etapa: "analisis"}).

\subsection{Verificación}

Después de escribir, el script comprueba dos cosas para cada fuente y umbral:
que haya 170 documentos y que la suma de \campo{n\_edificios} sea igual al
número de edificios asignados que pasan el filtro de confianza. Si alguna
falla, el script termina con error. Las cuatro combinaciones pasaron:

\begin{itemize}
  \item Google 0,65: 2\,694\,785 edificios, que es el total de la colección
        (2\,694\,786) menos el único sin municipio.
  \item Google 0,70 y 0,75: 2\,263\,445 y 1\,751\,076, iguales al conteo
        directo de edificios con \campo{cod\_mpio} y confianza mayor o igual
        al umbral.
  \item Microsoft: 1\,761\,048, el total de la colección.
\end{itemize}

En Google, 5 municipios quedan en cero con cualquier umbral (Murindó, Vigía
del Fuerte, Bojayá, Medio San Juan y Nóvita, los mismos del EDA) y Sipí pasa
a cero desde 0,70 porque su único edificio tiene confianza menor a 0,70. En
Microsoft los 170 municipios tienen edificios.

\subsection{Resultados por subregión}

\campo{exportar.py} agrupa \campo{resultados} por subregión con
\campo{\$lookup} y \campo{\$group} en MongoDB y genera
\ruta{data/processed/resultados_municipio.csv} (680 filas: 170 municipios por
cuatro combinaciones de fuente y umbral) y
\ruta{data/processed/resultados_subregion.csv}. Los dos CSV se versionan para
que los mapas y el ranking se puedan rehacer sin la base.

\begin{table}[H]
\centering\small
\caption{Edificios y área de techos por subregión PDET (Google con umbral 0,65). $D = (A_G - A_M)/(A_G + A_M)$; positivo si Google tiene más área.}
\label{tab:subregion}
\begin{tabularx}{\textwidth}{@{}>{\raggedright\arraybackslash}Xrrrrrr@{}}
\toprule
 & & \multicolumn{2}{c}{\textbf{Edificios}} & \multicolumn{2}{c}{\textbf{Área de techos (km\textsuperscript{2})}} & \\
\cmidrule(lr){3-4}\cmidrule(lr){5-6}
\textbf{Subregión} & \textbf{Mpios.} & \textbf{Google} & \textbf{Microsoft} & \textbf{Google} & \textbf{Microsoft} & $\boldsymbol{D}$ \\
\midrule
Sierra Nevada - Perijá & 15 & 645\,009 & 295\,383 & 53,46 & 46,27 & $+$0,07 \\
Alto Patía y Norte del Cauca & 24 & 388\,378 & 271\,668 & 30,86 & 30,36 & $+$0,01 \\
Cuenca del Caguán y Piedemonte Caqueteño & 17 & 182\,694 & 133\,941 & 20,93 & 22,57 & $-$0,04 \\
Urabá Antioqueño & 8 & 218\,242 & 113\,400 & 18,19 & 15,34 & $+$0,09 \\
Montes de María & 15 & 215\,874 & 129\,388 & 14,23 & 11,82 & $+$0,09 \\
Putumayo & 9 & 136\,711 & 108\,450 & 12,58 & 13,86 & $-$0,05 \\
Bajo Cauca y Nordeste Antioqueño & 13 & 142\,783 & 102\,406 & 10,91 & 10,69 & $+$0,01 \\
Macarena - Guaviare & 12 & 113\,455 & 94\,110 & 10,17 & 11,49 & $-$0,06 \\
Arauca & 4 & 111\,739 & 74\,675 & 9,42 & 10,20 & $-$0,04 \\
Sur de Córdoba & 5 & 114\,887 & 80\,634 & 8,70 & 8,33 & $+$0,02 \\
Pacífico Medio & 4 & 116\,255 & 49\,596 & 8,65 & 8,78 & $-$0,01 \\
Catatumbo & 8 & 104\,278 & 89\,633 & 8,21 & 9,22 & $-$0,06 \\
Sur de Bolívar & 7 & 80\,076 & 52\,856 & 6,03 & 5,84 & $+$0,02 \\
Pacífico y Frontera Nariñense & 11 & 71\,020 & 79\,144 & 5,30 & 8,17 & $-$0,21 \\
Sur del Tolima & 4 & 32\,662 & 48\,072 & 2,80 & 4,77 & $-$0,26 \\
Chocó & 14 & 20\,722 & 37\,692 & 1,73 & 3,43 & $-$0,33 \\
\midrule
\textbf{Total} & \textbf{170} & \textbf{2\,694\,785} & \textbf{1\,761\,048} & \textbf{222,17} & \textbf{221,13} & \textbf{0,00} \\
\bottomrule
\end{tabularx}
\end{table}


Las dos fuentes casi coinciden en área total (222,17 y 221,13\,km\textsuperscript{2})
aunque Google tiene 1,53 veces más edificios, porque divide la misma
superficie en huellas más pequeñas. Por subregión, Google tiene más área en el
Caribe (Sierra Nevada - Perijá, Urabá, Montes de María) y Microsoft en el
Pacífico y el sur (Chocó, Sur del Tolima, Pacífico y Frontera Nariñense), donde
$D$ llega a $-0{,}33$. Los conteos por subregión coinciden con los del EDA de
la Entrega~3 salvo en Catatumbo, con un edificio de diferencia: el EDA usó el
centroide que trae Google y aquí se usa el \campo{centroide} recalculado en
EPSG:9377.

\subsection{Sensibilidad al umbral de confianza}

\campo{sensibilidad.py} compara los tres umbrales de Google con Microsoft
como referencia.

\begin{table}[H]
\centering\small
\caption{Sensibilidad al umbral de confianza de Google. Las pérdidas son frente a 0,65. \emph{Mpios.}: municipios con al menos un edificio. \emph{Top 15}: cuántos de los 15 municipios con más área de techos con 0,65 siguen en el top 15.}
\label{tab:sensibilidad}
\begin{tabular}{@{}lrrrrrr@{}}
\toprule
 & \multicolumn{2}{c}{\textbf{Edificios}} & \multicolumn{2}{c}{\textbf{Área (km\textsuperscript{2})}} & & \\
\cmidrule(lr){2-3}\cmidrule(lr){4-5}
\textbf{Escenario} & \textbf{Total} & \textbf{Pérdida} & \textbf{Total} & \textbf{Pérdida} & \textbf{Mpios.} & \textbf{Top 15} \\
\midrule
Google, conf.\ $\geq$ 0,65 & 2\,694\,785 & -- & 222,17 & -- & 165 & 15/15 \\
Google, conf.\ $\geq$ 0,70 & 2\,263\,445 & 16,0\,\% & 206,67 & 7,0\,\% & 164 & 15/15 \\
Google, conf.\ $\geq$ 0,75 & 1\,751\,076 & 35,0\,\% & 179,63 & 19,1\,\% & 164 & 14/15 \\
\midrule
Microsoft (referencia) & 1\,761\,048 &  & 221,13 &  & 170 & 11/15 \\
\bottomrule
\end{tabular}
\end{table}


Subir el umbral quita muchos más edificios que área: con 0,70 se pierde el
16\,\% de los edificios pero solo el 7\,\% del área, y con 0,75 el 35\,\% y el
19\,\%. Los edificios de baja confianza son, sobre todo, huellas pequeñas.
Aun así, el área con 0,75 (179,6\,km\textsuperscript{2}) queda 19\,\% por
debajo de Microsoft, mientras que con 0,65 las dos fuentes casi coinciden.
El orden de los municipios es estable: los 15 con más área de techos con 0,65
siguen en el top 15 con 0,70, y 14 de 15 con 0,75. Por eso se deja 0,65 como
escenario base, y la sensibilidad muestra que la elección del umbral afecta
el total pero no cuáles municipios encabezan el ranking.
